% =====================================================================
% Chapter 5, part 3: Future work. Transcribed from the reviewed
% DISCUSSION.html draft (section 7).
% =====================================================================

\section{Future Work}
\label{sec:futurework}

The limitations of Section~\ref{sec:limitations} are not dead ends. Ordered
from immediately actionable to long-term, Table~\ref{tab:futurework} lays
out how stronger and trustworthy results can be built on this line of work.

\begin{table}[H]
  \centering
  \caption[Future work: the path to stronger results]{Future work: the path to stronger results.}
  \label{tab:futurework}
  \small
  \begin{tabular}{@{}L{4.4cm}L{9.9cm}@{}}
    \toprule
    \textbf{Step} & \textbf{What it would deliver} \\
    \midrule
    \textbf{1. Model the timing signature directly.} &
    Extract turn duration, sit-to-stand time and cadence from the
    micro-Doppler signature as explicit quantities, instead of hoping a network
    discovers them in pixels. Timing is the one signal this work showed to
    be real: patients take longer, and the extra time concentrates in
    standing up and turning. A model built directly on these timing
    quantities would be an automated version of the timed-up-and-go test
    that clinicians already use, which is also what the source radar system
    was validated to measure \cite{lopezdelgado2026}. \\
    \addlinespace
    \textbf{2. Obtain subject ages.} &
    An age-matched analysis would settle whether the surviving signal is
    disease or ageing, and could also unmask real signal that age variation
    currently hides. \\
    \addlinespace
    \textbf{3. Record a larger, session-balanced cohort.} &
    With 58 subjects every score carries a wide uncertainty range, so even a
    genuinely better model could not prove itself here. Demonstrating a
    spectral gait signature clearly above the duration floor needs a cohort
    in the low hundreds, recorded in sessions balanced across the two
    groups. \\
    \addlinespace
    \textbf{4. Keep the evaluation standard.} &
    Whatever the model, evaluate it on people it has never seen and report
    it next to the duration floor and a confidence interval. The protocol
    built for this thesis is pre-registered, ablation-tested and reusable as
    the design document for every step above. \\
    \bottomrule
  \end{tabular}
\end{table}

The review of this work also pointed to methodological improvements that
target the specific failure found here, a network that looks at the wrong
part of the image. Table~\ref{tab:futuremethods} lists them, each with the
work that introduced the idea.

\begin{table}[H]
  \centering
  \caption[Future work: methodological improvements]{Future work: methodological improvements suggested at review.}
  \label{tab:futuremethods}
  \small
  \begin{tabular}{@{}L{4.4cm}L{9.9cm}@{}}
    \toprule
    \textbf{Technique} & \textbf{What it would change} \\
    \midrule
    \textbf{Make the network look at the gait.} &
    Attention-guided masking trains the attention map itself to stay inside
    the region of interest \cite{li2018gain}; explanation regularisation
    penalises attention on regions known to be irrelevant, so that the network
    is right for the right reasons \cite{ross2017,rieger2020}; counterfactual
    training adds examples in which only the irrelevant part of the image
    changes, so that it cannot carry the label \cite{kaushik2020}. All three
    act directly on the failure the attention maps revealed. \\
    \addlinespace
    \textbf{Remove the background first.} &
    An adaptive threshold set from the local energy of each signature strips
    the background before training and leaves only the moving-body content
    \cite{li2023adaptive}. Combined with the masking above, it takes the
    shortcut away instead of asking the network to ignore it. \\
    \addlinespace
    \textbf{Learn a compact embedding first.} &
    An encoder with a narrow bottleneck compresses each signature into a few
    informative numbers before any classifier sees it
    \cite{seyfioglu2018dcae}. Plotting that embedding with t-SNE or UMAP
    \cite{vandermaaten2008,mcinnes2018} shows at a glance whether it separates
    by diagnosis, by person or by recording length. \\
    \addlinespace
    \textbf{Borrow encoders from speech.} &
    A micro-Doppler signature is closer to a speech spectrogram than to a
    photograph, and encoders pretrained on speech transfer to it better than
    image networks do \cite{li2020speech,tran2020}; the idea goes back to
    speech-recognition front-ends applied to micro-Doppler \cite{yessad2011}.
    This replaces the ImageNet probe with a source domain that matches. \\
    \addlinespace
    \textbf{Estimate the Doppler content parametrically.} &
    An autoregressive spectral estimate replaces the short-time Fourier
    transform with a model fitted to each short segment, giving smoother
    signatures with less noise from the same data \cite{kaymarple1981}. \\
    \addlinespace
    \textbf{Tune the window length.} &
    The 3.0\,s window was fixed in advance from the stride duration. Treating
    it as a hyperparameter, chosen with the AUC as figure of merit inside the
    inner validation loop, would test whether a longer or shorter view helps. \\
    \bottomrule
  \end{tabular}
\end{table}

\section{Closing Remarks}
\label{sec:closing}

This work asked a simple clinical question of a rich dataset and refused the
shortcuts that would have made the answer look better than it is. What it
hands to the next study is a modest and honest number, a clear explanation of
why that number is not higher, and an evaluation protocol that any future
attempt can be held to. The day richer data arrives, with ages, severity
scores and the original per-sensor recordings, the standard defended here
will be ready to measure it, on people the model has never seen.
